\chapter{Additional Materials}\label{chapter:appendix}

Bu bölümde tezde yer almayan ek materyaller sunulmaktadır.

\section{Configuration Parameters Reference}

Table~\ref{tab:cfg_full} provides a complete listing of all MycoFlow UCI configuration parameters, their types, default values, and acceptable ranges.

\begin{table}[h]
\caption{Complete MycoFlow UCI Configuration Reference}
\label{tab:cfg_full}
\centering
\small
{\setlength{\tabcolsep}{4pt}
\begin{tabular}{p{4cm}p{1.5cm}p{2cm}p{5cm}}
\toprule
\textbf{Parameter} & \textbf{Type} & \textbf{Default} & \textbf{Description} \\
\midrule
\texttt{sample\_hz}         & float & 1.0     & Sense cycle frequency; range [0.1, 10.0] Hz \\
\texttt{ewma\_alpha}        & float & 0.30    & EWMA smoothing factor; range (0, 1) \\
\texttt{rtt\_margin}        & float & 0.30    & Congestion threshold margin $\gamma$; range (0, 1) \\
\texttt{baseline\_decay}    & float & 0.01    & Baseline drift rate $\beta$; range (0, 0.1) \\
\texttt{baseline\_interval} & int   & 60      & Cycles between baseline updates \\
\texttt{bw\_step}           & int   & 2000    & Bandwidth adjustment step (kbit/s) \\
\texttt{bw\_min}            & int   & 1000    & Minimum WAN cap (kbit/s) \\
\texttt{bw\_max}            & int   & 1000000 & Maximum WAN cap (kbit/s) \\
\texttt{action\_cooldown}   & float & 3.0     & Minimum action interval (s) \\
\texttt{baseline\_samples}  & int   & 5       & Startup calibration sample count \\
\texttt{per\_device}        & bool  & 1       & Enable per-device DSCP marking \\
\texttt{flow\_aware}        & bool  & 1       & Enable v3 per-flow classifier + CONNMARK \\
\texttt{ingress\_enabled}   & bool  & 0       & Enable IFB ingress shaping \\
\texttt{wan\_interface}     & string & eth0   & WAN interface name \\
\texttt{probe\_host}        & string & 1.1.1.1 & ICMP probe target host \\
\texttt{rtt\_bpf\_obj}      & string & ---    & Path to mycoflow\_rtt.bpf.o \\
\bottomrule
\end{tabular}}
\end{table}

\section{Unicauca Dataset Service Class Mapping}

Table~\ref{tab:unicauca_map} shows the mapping between Unicauca dataset application labels and MycoFlow service classes used in the Phase~4 evaluation.

\begin{table}[h]
\caption{Unicauca Application Label to MycoFlow Service Class Mapping}
\label{tab:unicauca_map}
\centering
\small
\begin{tabular}{ll}
\toprule
\textbf{Unicauca Label} & \textbf{MycoFlow Service Class} \\
\midrule
Game (UDP)     & \texttt{game\_rt} \\
VoIP           & \texttt{voip\_call} \\
Video Call     & \texttt{video\_conf} \\
Live Streaming & \texttt{video\_live} \\
Video on Demand & \texttt{video\_vod} \\
Web Browsing   & \texttt{web\_interactive} \\
File Transfer  & \texttt{file\_sync} \\
P2P (BitTorrent) & \texttt{torrent} \\
FTP / Bulk     & \texttt{bulk\_dl} \\
Game Launcher  & \texttt{game\_launcher} \\
DNS / System   & \texttt{system} \\
\bottomrule
\end{tabular}
\end{table}

\section{CAKE Tin Target Delay Summary}

CAKE's internal AQM target delay is computed from the \texttt{rtt} parameter as $\text{target} \approx \text{rtt}/20$. Table~\ref{tab:cake_targets} shows the AQM targets applied per persona:

\begin{table}[h]
\caption{Persona-Driven CAKE AQM Target Delays}
\label{tab:cake_targets}
\centering
\small
\begin{tabular}{lcc}
\toprule
\textbf{Persona Group} & \textbf{CAKE rtt param} & \textbf{AQM target} \\
\midrule
VOIP, GAMING    & 50~ms  & 2.5~ms \\
VIDEO           & 75~ms  & 3.75~ms \\
STREAMING       & 150~ms & 7.5~ms \\
BULK, TORRENT   & 200~ms & 10~ms \\
UNKNOWN         & 100~ms & 5~ms (CAKE default) \\
\bottomrule
\end{tabular}
\end{table}

\section{Build and Installation Instructions}

\subsection{Building from Source}

\begin{lstlisting}[language=bash, basicstyle=\ttfamily\small]
# Install OpenWrt SDK for aarch64-openwrt-linux-musl
# Set SDK_PATH to the SDK root directory

mkdir build && cd build
cmake .. -DCMAKE_TOOLCHAIN_FILE=${SDK_PATH}/toolchain.cmake \
         -DCMAKE_BUILD_TYPE=Release
make -j$(nproc)
ctest --output-on-failure  # Run all 13 test suites
\end{lstlisting}

\subsection{OpenWrt Installation}

\begin{lstlisting}[language=bash, basicstyle=\ttfamily\small]
# Copy binary to router
scp mycoflowd root@192.168.1.1:/usr/sbin/

# Install required kernel module (OpenWrt 24.10+)
opkg install kmod-nf-conntrack-netlink

# Configure via UCI
uci set mycoflow.main.wan_interface=eth1
uci set mycoflow.main.flow_aware=1
uci commit mycoflow

# Start daemon
/etc/init.d/mycoflow start
\end{lstlisting}
